\color{uniblue}\section[Выбор уставок дифференциальной токовой защиты трансформатора]{\sloppy Выбор уставок дифференциальной токовой защиты \mbox{трансформатора}} 
\color{black}

\begin{enumerate}[label=\arabic{section}.\arabic{enumi}, labelsep=7pt, leftmargin=0pt, itemindent=60pt, align=left, nosep]
\item
Функция ДЗТ предназначена для быстрого и селективного отключения объекта защиты при внутренних повреждениях всех видов. Для каждой фазы предусмотрены дифференциальный токовый орган с торможением (ДТЗт) и орган дифференциальной токовой отсечки (ДТО). Состав функции обеспечивает цифровую компенсацию фазного сдвига и выравнивание подводимых токов сторон (плеч). Также обеспечивается отстройка от режима броска намагничивающего тока и режима перевозбуждения трансформатора. Для надежной работы функции при высоком уровне токов внешних КЗ, вызывающих насыщение трансформаторов тока может быть введен дополнительный критерий внешних повреждений на основе сравнения фаз токов по сторонам защищаемого трансформатора. В составе функции ДЗТ также предусмотрен контроль исправности цепей тока для каждой из сторон.

\item \textbf{Выбор общих уставок}

\begin{enumerate}[label=\arabic{section}.\arabic{enumi}.\arabic*, labelsep=7pt, leftmargin=0em, itemindent=55.5pt, align=left, nosep]

\item
Функция ДЗТ в устройстве ЮНИТ-М300-ДЗТ2 позволяет подключить до трех плеч защит. В зависимости от  количества обмоток защищаемого силового трансформатора необходимо с помощью уставок определить используемые аналоговые входы переменного тока устройства. Например, параметр <<Сторона 3>>, его возможные значения -- <<Не используется>>, <<Используется>>.
Далее для используемых аналоговых входов сторон определить схему соединения вторичных обмоток измерительных трансформаторов тока (<<звезда>> или <<треугольник>>).
Необходимо также задать значение базисной мощности, которое допускается принять равным номинальной мощности трансформатора.

\item Дополнительные параметры приведения сторон (задаются для каждой из используемых сторон). Параметр <<Номинальное напряжение стороны n>> – должен соответствовать паспортным значениям номинального напряжения стороны n трансформатора.
Параметры <<Номинальный первичный ток трансформатора тока стороны n>> и <<Номинальный вторичный ток трансформатора тока стороны n>> также должны соответствовать номинальным параметрам используемых в данном плече защиты трансформаторов тока.
\item
По указанным параметрам (<<Номинальное напряжение стороны n>>, <<Номинальный первичный ток трансформатора тока стороны n>> и <<Номинальный вторичный ток трансформатора тока стороны n>>) устройство автоматически определяет коэффициент амплитудного выравнивания для рассматриваемой стороны. Рассмотренные ниже параметры (<<Компенсация токов 3I0 для стороны n>> и <<Векторная группа стороны n>>) используются для определения устройством коэффициентов матрицы компенсации фазового сдвига.
\item
Параметр <<Компенсация токов 3I0 для стороны n>>. Компенсацию токов 3I0 рекомендуется задействовать на стороне силового трансформатора, нейтраль обмотки которой заземлена, для исключения возможности ложного срабатывания ДЗТ при внешних КЗ на землю. При соединении вторичных обмоток ТТ стороны в <<треугольник>> компенсацию токов 3I0 вводить не следует.
\item
Параметр <<Векторная группа стороны n>> предназначен для компенсации фазового сдвига между токами, измеряемыми аналоговыми входами разных сторон, обусловленного схемой соединения  обмоток защищаемого трансформатора (а также возможным различием в схеме соединения вторичных обмоток ТТ по сторонам трансформатора). Группа соединения, умноженная на 30°, определяет на какой электрический угол по направлению часовой стрелки должен быть сдвинут вектор фазного тока для корректного сравнения. Характерные значения – 0, 1, 5, 6, 11.
\end{enumerate} 

\item \textbf{Выбор уставок дифференциального токового органа с торможением}

\begin{enumerate}[label=\arabic{section}.\arabic{enumi}.\arabic*, labelsep=7pt, leftmargin=0em, itemindent=55.5pt, align=left, nosep]
\item
Характеристика срабатывания ДТЗт представляет из себя ломаную линию в координатах дифференциального тока и тока торможения (см. рисунок \ref{dzt:img1}). Данная характеристика (выделена синим цветом) состоит из трех участков (начальный, первый и второй наклонные), определяется следующими параметрами:
\begin{itemize} 
\item дифференциальный ток срабатывания начального участка, $I_\textsl{ср(ДТЗт)}$;
\item значение тока торможения, соответствующее началу первого наклонного участка, $I_\textsl{т1}$;
\item коэффициент торможения первого наклонного участка, $K_\textsl{т1}$;
\item значение тока торможения, соответствующее началу второго наклонного участка, $I_\textsl{т2}$;
\item коэффициент торможения второго наклонного участка, $K_\textsl{т2}$.
\end{itemize}
При неисправностях в токовых цепях характеристика загрубляется на начальном участке (выделена зеленым цветом).
Следует отметить, что параметры тормозной характеристики, заданные по умолчанию подходят в большинстве случаев применения ДТЗт.
\vspace{3mm}
\begin{figure}[H]
\centering
\includegraphics[width=0.7\textwidth,height=0.7\textheight,keepaspectratio]{\fbpath/08. ДЗТ/dzt.pdf}
\captionof{figure}{Характеристика срабатывания дифференциального токового органа с торможением и дифференциальной токовой отсечки}\label{dzt:img1}
\end{figure}

\item 
Согласно \cite{ru13b} токи небаланса (из допущения, что расчет ведется только для периодической составляющей) в общем случае определяются как сумма трех составляющих, которые обусловлены погрешностями измерения, устройством регулирования напряжения и выравниванием токов плеч.

\begin{equation}
    \label{eq:dzt1}
    I_{\textsl{нб.расч.}} =  I_{\textsl{нб.расч.*}}^{'} + I_{\textsl{нб.расч.*}}^{''} + I_{\textsl{нб.расч.*}}^{'''} ,
\end{equation}
\begin{equation}
    \label{eq:dzt2}
    I_{\textsl{нб.расч.*}}^{'}=  k_{\textsl{пер}} \cdot k_{\textsl{одн}} \cdot \epsilon \cdot I_{\textsl{расч}} ,
\end{equation}
\begin{equation}
    \label{eq:dzt3}
    I_{\textsl{нб.расч.*}}^{''}= \Delta U_{\textsl{рег}} \cdot I_{\textsl{расч}} ,
\end{equation}
\begin{equation}
    \label{eq:dzt4}
    I_{\textsl{нб.расч.*}}^{'''}= f_{\textsl{выр}} \cdot I_{\textsl{расч}} ,
\end{equation}
\begin{eqexpl}[15mm]
    \item {$I_{\textsl{нб.расч.*}}^{'}$} составляющая тока небаланса, обусловленная погрешностью ТТ;
    \item{$k_{\textsl{пер}}$} коэффициент переходного режима (учитывает наличие апериодической составляющей);
    \item{$k_{\textsl{одн}}$} коэффициент однотипности трансформаторов тока (рекомендуется принимать с запасом 1,0);
    \item{$\epsilon$} относительное значение полной погрешности трансформаторов тока (всегда принимается равным 0,1);
    \item {$I_{\textsl{нб.расч.*}}^{''}$} составляющая тока небаланса, обусловленная регулированием напряжения защищаемого трансформатора, при отсутствии регулирования данная составляющая не принимается в расчет, а при постоянной работе трансформатора на конкретной отпайке, необходимо параметр <<Номинальное напряжение стороны n>> (n – в данном случае сторона, где установлено устройство регулирования) задавать соответствующим данной отпайке;
    \item{$\Delta U_{\textsl{рег}}$} погрешность, обусловленная регулированием напряжения защищаемого трансформатора, принимается равной максимальному отклонению напряжения от номинального в относительных величинах;
    \item {$I_{\textsl{нб.расч.*}}^{'''}$} составляющая тока небаланса, обусловленная погрешностью выравнивания токов плеч;
    \item{$f_{\textsl{выр}}$} погрешность выравнивания токов плеч, обусловленное погрешностями измерений; рекомендуется принимать равной 0,03;  
    \item{$I_{\textsl{расч}}$} относительный ток в режиме, для которого ведется расчет.   
\end{eqexpl}

\item 
Начальный участок характеристики определяется минимальным для всей характеристики дифференциальным током срабатывания, который отстраивается от токов небаланса в нормальном режиме работы трансформатора. Этот участок гарантирует выявление витковых замыканий в обмотках (когда ток повреждения может составлять меньше номинального).

\item 
Параметр <<Дифференциальный ток срабатывания начального участка>> выбирается по условию отстройки от тока небаланса в нормальном режиме по выражению:
\begin{equation}
    \label{eq:dzt5}
    I_{\textsl{ср}}= k_{\textsl{отс}} \cdot I_{\textsl{нб.расч}},
\end{equation}
\begin{eqexpl}[15mm]
    \item {$k_{\textsl{отс}}$} коэффициент отстройки, учитывающий погрешности аналоговых входов терминала, принимается равным 1,1$\div$1,3;
    \item {$I_{\textsl{нб.расч}}$} расчетный ток небаланса в нагрузочном режиме, определяется по формуле (\ref{eq:dzt1}), при этом $I_{\textsl{расч}}$ принимается равным номинальному току (1,0), а параметр $k_{\textsl{пер}}$ принимается равным 1,0. 
\end{eqexpl}

\item 
Параметр $I_{\textsl{ср}}$ не рекомендуется принимать равным менее 0,2, рекомендуемые значения находятся в диапазоне 0,3$\div$0,5.

%%%%%%%%%%%%%%%% ТОК ТОРМОЖЕНИЯ ПЕРВОГО НАКЛОННОГО УЧАСТКА
\item
Первый наклонный участок характеристики служит для отстройки от максимальных рабочих нагрузок, перегрузок, эксплуатационных переключений, а также внешних коротких замыканий. Защита загрубляется ровно настолько, чтобы не отключить неповрежденный трансформатор из-за работы РПН на крайних отпайках или при тяжелом пуске/самозапуске двигательной нагрузки. При принятом в устройстве способе формирования тока торможения рекомендуется значение $I_{\textsl{т1}}$ принимать равным 1,0.
%%%%%%%%%%%%%%%% К ТОРМОЖЕНИЯ ВТОРОГО НАКЛОННОГО УЧАСТКА
\item
Коэффициент торможения $k_{\textsl{т1}}$ первого наклонного участка рекомендуется определять по выражению:
\begin{equation}
    \label{eq:dzt6}
    k_{\textsl{т1}} = \frac {I_{\textsl{ср(т2)}} - I_{\textsl{ср}}}{I_{\textsl{т2}}-I_{\textsl{т1}}},
\end{equation}
\begin{eqexpl}[15mm]
    \item {$I_{\textsl{ср(т2)}}$} расчетный дифференциальный ток, определяемый по выражениям (\ref{eq:dzt5}) и (\ref{eq:dzt1}), при этом:
    \par\hangindent=2.5em\hangafter=1
    \raggedright
    - параметр $I_{\textsl{расч}}$ принимается равным значению тормозного тока при внешнем максимальном КЗ (как правило, трехфазным), т.к. в режиме внешнего (сквозного) КЗ алгоритм полусуммы модулей токов плеч сводится к полной величине транзитного тока ($I_{\textsl{торм}}=I_{\textsl{скв}}$), что исключает необходимость введения понижающих коэффициентов при определении сквозного тока;\\
    \centering
    $I_{\textsl{расч}} = I_{\textsl{торм}} = \frac {I_{\textsl{КЗ, max}} \cdot k_{\textsl{сх}} \cdot k_{\textsl{ампл}} \cdot I_{\textsl{втор. ТТ}}}{I_{\textsl{перв. ТТ}}}$,\\
    \raggedright
    где\\ - коэффициент амплитудного выравнивания стороны трансформатора, для которой указан ток КЗ расчетного режима, рассчитывается по выражению:\\
    \centering    
    $k_{\textsl{ампл}} = \frac {I_{\textsl{перв. ТТ}}}{I_{\textsl{втор. ТТ}} \cdot k_{\textsl{сх}} \cdot I_{\textsl{баз}}}$,\\
    \raggedright   
    - базисный ток трансформатора $I_{\textsl{баз}}$ определяется через базисную мощность $S_{\textsl{баз}}$ и номинальное напряжение стороны $U_{\textsl{ном}}$:\\
    \centering  
    $I_{\textsl{баз}} = \frac {S_{\textsl{баз}}}{\sqrt(3) \cdot  U_{\textsl{ном}}}$,\\
    \raggedright      
    - параметр $k_{\textsl{отс}}$ коэффициент отстройки, принимается равным 1,1$\div$1,3;\\
    - параметр $k_{\textsl{пер}}$ принимается равным 2,5, если доля двигательной нагрузки более 50\%, в остальных случаях принимается равным 1,5$\div$2,0;\\
    - $I_{\textsl{т1}}$ значение тока торможения, соответствующее началу первого наклонного участка, о.е.;\\   
    - $I_{\textsl{т2}}$ тормозной ток в расчетном режиме. принимается равным максимальному относительному току при внешнем (как правило, трехфазном) КЗ, о.е.\\   
\end{eqexpl}


%%%%%%%%%%%%%%%% ТОК ТОРМОЖЕНИЯ ВТОРОГО НАКЛОННОГО УЧАСТКА

\item
В режиме, соответствующем началу второго наклонного участка, характеристика отстраивается от нелинейной погрешности ТТ. Затухающая апериодическая (постоянная) составляющая тока мгновенно насыщает магнитопроводы трансформаторов тока. В состоянии насыщения ТТ перестают корректно трансформировать первичный ток во вторичную цепь. Возникает дифференциальный ток (переходный ток небаланса), который терминал защиты может ложно принять за внутреннее повреждение. Значение $I_{\textsl{т2}}$ рекомендуется выбирать из диапазона 2,0$\div$3,0. 

\item
Коэффициент торможения $k_{\textsl{т2}}$ второго наклонного участка рекомендуется определять по выражению:
\begin{equation}
    \label{eq:dzt7}
    k_{\textsl{т2}} = \frac {I_{\textsl{ср.макс}} - I_{\textsl{ср(т2)}}}{I_{\textsl{т2макс}}-I_{\textsl{т2}}},
\end{equation}
\begin{eqexpl}[15mm]
    \item {$I_{\textsl{ср.макс}}$} расчетный дифференциальный ток, определяемый по выражениям (\ref{eq:dzt5}) и (\ref{eq:dzt1}), при этом:
    \par\hangindent=2.5em\hangafter=1    
     - параметр $I_{\textsl{расч}}$ принимается равным коэффициенту предельной кратности ТТ при реальной нагрузке $k'_{10}$ (при отсутствии данных допустимо принять равным 20);\\ 
     - параметр $k_{\textsl{пер}}$ принимается равным 3,0$\div$4,0;
    \item {$I_{\textsl{ср(т2)}}$} относительный расчетный дифференциальный ток, соответствующий концу первого участка тормозной характеристики, о.е, рассчитывается по выражению: \\{$I_{\textsl{ср(т2)}} = I_{\textsl{ср(т1)}}+k_{\textsl{т1}}(I_{\textsl{т1}}-I_{\textsl{т2}})$};

    \item {$I_{\textsl{т2макс}}$} тормозной ток расчетного режима, принимается равным коэффициенту предельной кратности ТТ при реальной нагрузке $k'_{10}$ (при отсутствии данных допустимо принять равным 20);

    \item {$I_{\textsl{т2}}$} ток торможения, соответствующий началу второго наклонного участка.
\end{eqexpl}

\item
Для обеспечения селективности ДТЗ полная погрешность измерительных трансформаторов тока в установившемся режиме максимального сквозного тока не должна превышать 10\%. Выбор ТТ базируется на параметре приведенной предельной кратности $k'_{10}$. 

\item
\needspace{3\baselineskip}
Приведенная предельная кратность вычисляется по выражению:
\begin{equation}
    \label{eq:dzt71}
    k'_{10} = \frac {I_{\textsl{перв}}}{I_{\textsl{ном}}}\cdot k_{10},
\end{equation}
\begin{eqexpl}[15mm]
\item  {$I_{\textsl{перв}}$} первичный номинальный ток ТТ стороны N, А;
\item  {$I_{\textsl{ном}}$} номинальный ток защищаемого трансформатора стороны N, А;
\item  {$k_{10}$} предельная кратность ИТТ при номинальной нагрузке\footnote{~Предельная кратность при номинальной нагрузке ($k_{10}$) -- это паспортная характеристика ТТ класса защиты (например, 10Р), которая показывает, во сколько раз первичный ток может превысить номинальный первичный ток самого ИТТ ($I_{\textsl{перв}}$), прежде чем полная погрешность трансформации в установившемся режиме превысит допустимые 10\%.}.
\end{eqexpl}

\item
Значение коэффициента торможения второго наклонного участка должно соответствовать условию:
\begin{equation}
    \label{eq:dzt8}
    k_{\textsl{т1}} \leq k_{\textsl{т2}} \leq 1.
\end{equation}

\item 
Согласно \cite{ru13b} чувствительность защиты определяется при металлическом КЗ на выводах защищаемого трансформатора при его работе на расчетном ответвлении. Расчетными режимами работы подстанции и питающих систем являются реальные режимы, обуславливающие минимальный ток при расчетном виде КЗ. В соответствие с \cite{prikaz546} наименьший коэффициент чувствительности должен составлять 2,0. В отдельных случаях, согласно \cite{prikaz546} значение требуемого коэффициента чувствительности может быть снижено до 1,5.


\item Проверка чувствительности ДТЗ производится при минимальном внутреннем токе короткого замыкания. Оценка базируется на аналитическом методе, определяющем фактический геометрический запас до границы срабатывания с учетом пропорционального снижения дифференциального и тормозного токов при развитии аварии.

\item В связи с аппаратным алгоритмом терминала, формирующим тормозной сигнал как полусумму фазных токов, при одностороннем внутреннем повреждении принимается следующее соотношение:

Значение коэффициента торможения второго наклонного участка должно соответствовать условию:
\begin{equation}
    \label{eq:dzt81}
    I_{\textsl{дифф}} = I_{\textsl{КЗ,мин}},
\end{equation}
\begin{equation}
    \label{eq:dzt82}
    I_{\textsl{торм}} = 0,5 \cdot I_{\textsl{КЗ,мин}},
\end{equation}

%%%%%%%%%%%%%%%%%%%%%%%%%%%%%%%%%%%%%%%%%%%%%%%%%%%%%%%%%%%%%% Коэффициент чувствительности 

\item \textbf{Определение коэффициента чувствительности}

\item
Согласно \cite{ru13b}, чувствительность защиты определяется при металлическом КЗ на выводах защищаемого трансформатора при его работе на расчетном ответвлении. Расчетными режимами работы подстанции и питающих систем являются реальные режимы, обуславливающие минимальный ток при расчетном виде КЗ. В соответствии с \cite{pue} наименьший коэффициент чувствительности должен составлять 2,0. В отдельных случаях, согласно \cite{pue}, значение требуемого коэффициента чувствительности может быть снижено до 1,5.

\item
Для подтверждения выполнения данных требований применяется аналитический метод, оценивающий чувствительность защиты динамически. Коэффициент чувствительности вычисляется как соотношение длин отрезков на прямой, описывающей траекторию перехода трансформатора из предшествующего нормального нагрузочного режима в аварийный режим короткого замыкания.

\item
Расчет базируется на алгоритме, в котором тормозной ток формируется через полусумму токов плеч. В режиме внутреннего КЗ с односторонним питанием расчетный тормозной ток равен половине дифференциального:  $I_{\textsl{торм}} = 0,5 \cdot I_{\textsl{дифф}}$. Данный коэффициент (0,5) является базовым для приведения формул.


\item  \textbf{Порядок расчета коэффициента чувствительности}

\item
Требуется определить расчетный дифференциальный ток $I_{\textsl{дифф}}$ при минимальном токе КЗ в конце защищаемой зоны в режиме одностороннего питания.

\begin{equation}
    \label{eq:dzt91}
    I_{\textsl{дифф}} = \frac {I_{\textsl{КЗ,мин}} \cdot k_{\textsl{сх}}\cdot  \cdot k_{\textsl{ампл}}\cdot I_{\textsl{втор.ТТ}} }{I_{\textsl{перв.ТТ}}},
\end{equation}


\item  \textbf{Расчет по полной аналитической формуле}
\item
С учетом траектории перехода из нагрузочного режима и коэффициента формирования тормозного тока расчетное выражение имеет вид:
\begin{equation}
    \label{eq:dzt92}
    k_{\textsl{ч}} = \frac {I_{\textsl{дифф}} - k_{\textsl{т1}}\cdot (0,5\cdot I_{\textsl{дифф}} - I_{\textsl{нагр*}}) }{I_{\textsl{ср(ДТЗт)}}+ k_{\textsl{т1}} \cdot (I_{\textsl{нагр*}} - I_{\textsl{т1}}) },
\end{equation}
\begin{eqexpl}[15mm]
\item  {$I_{\textsl{нагр*}}$} относительный ток нагрузки предшествующего режима, о.е.;
\item  {$I_{\textsl{дифф}}$} дифференциальный ток КЗ, который возникает при минимальном внутреннем КЗ.
\end{eqexpl}

\item  \textbf{Расчет по упрощенной аналитической формуле}
\item
Допускается применение допущения, при котором относительный ток нагрузки трансформатора до момента КЗ равен току излома характеристики: $I_{\textsl{дифф}}=I_{\textsl{торм}}=1,0$ о.е. Выражение принимает упрощенный вид: 
\begin{equation}
    \label{eq:dzt92}
    k_{\textsl{ч}} = \frac {I_{\textsl{дифф}} - k_{\textsl{т1}}\cdot (0,5\cdot I_{\textsl{дифф}} - 1) }{I_{\textsl{ср(ДТЗт)}} },
\end{equation}
\item
Для обеспечения надежности срабатывания коэффициент должен удовлетворять условию $k_{\textsl{ч}} \geq 2,0$.


\item  \textbf{Проверка корректности расчета}
\item 
Аналитический метод выдает достоверный результат исключительно в том случае, если расчетная прямая, соединяющая точку нагрузки с точкой КЗ, пересекает характеристику срабатывания строго на ее первом наклонном участке. Проверка ординаты точки пересечения выполняется в два этапа:

\begin{itemize}
\item Вычисление ординаты пересечения расчетной прямой с тормозной характеристикой:
\begin{equation}
    \label{eq:dzt921}
    I_{\textsl{дифф, тх*}} = \frac {I_{\textsl{дифф}} }{k_{\textsl{ч}} },
\end{equation}
\item Контроль граничного условия участка. Полученное значение должно быть меньше или равно верхней границе первого наклонного участка:
\begin{equation}
    \label{eq:dzt922}
    I_{\textsl{дифф, тх*}} \leq I_{\textsl{ср(ДТЗт)}} + k_{\textsl{т1} \cdot (I_{\textsl{т2}}-I_{\textsl{т1}})}, 
\end{equation}
Если неравенство не выполняется (ордината пересечения выше границы участка), использование данной аналитической формулы неправомерно, и требуется корректировка уставок излома характеристики.

\end{itemize}






\end{enumerate}

\item \textbf{Выбор уставок функций блокировок по второй и пятой гармоникам дифференциального тока}

\begin{enumerate}[label=\arabic{section}.\arabic{enumi}.\arabic*, labelsep=7pt, leftmargin=0em, itemindent=55.5pt, align=left, nosep]
\item
На основании опыта производителей релейной защиты соотношение токов второй и первой гармоник на уровне 12–15\% обеспечивают надежную блокировку дифференциальной токовой защиты с торможением при бросках намагничивающего тока.
\item
В сетях 35 кВ перевозбуждение силового трансформатора маловероятно, поэтому блокировку по наличию пятой гармоники в дифференциальном токе допускается не вводить в работу.
\end{enumerate}

\item Выбор уставок дифференциальной токовой отсечки
\begin{enumerate}[label=\arabic{section}.\arabic{enumi}.\arabic*, labelsep=7pt, leftmargin=0em, itemindent=55.5pt, align=left, nosep]
\item
Дифференциальная токовая отсечка предназначена для мгновенного отключения повреждений в зоне действия дифференциальной защиты, сопровождаемых значительными токами КЗ. ДТО является более грубым, чем ДТЗт органом и должна быть отстроена от бросков тока намагничивания и небаланса при внешних КЗ. На рисунке \ref{dzt:img1} характеристика срабатывания ДТО представлена красной линией.
Отстройка от бросков тока намагничивания обеспечивается в общем случае (согласно \cite{ru13b}) при выполнении условия:
\begin{equation}
    \label{eq:dzt16}
    I_{\textsl{ср(ДТО)}} \geq 6 ,
\end{equation}
Следует иметь ввиду, что чем меньше мощность трансформатора, тем выше значение кратности тока намагничивания.
Отстройка от максимального тока небаланса определяется по (\ref{eq:dzt5}) и (\ref{eq:dzt1}), при этом:

\begin{itemize} %[label=-, labelwidth=-0.5em, labelsep=0.5em, align=left, leftmargin=0em, itemindent=2.5em, nosep]
\item параметр $I_{\textsl{расч}}$ принимается равным максимальному значению тока внешнего КЗ (как правило трехфазного);
\item параметр $k_{\textsl{отс}}$ принимается равным 1,5;
\item параметр $k_{\textsl{пер}}$ принимается равным 3,0$\div$4,0.
\end{itemize}
Из полученных по указанным выше условиям значений $I_{\textsl{ср(ДТО)}}$ в качестве уставки выбирается большее.
\end{enumerate}

\item Выбор уставок функции контроля цепей тока
\begin{enumerate}[label=\arabic{section}.\arabic{enumi}.\arabic*, labelsep=7pt, leftmargin=0em, itemindent=55.5pt, align=left, nosep]
\item
В устройстве ЮНИТ-М300-ДЗТ2 существует два способа контроля исправности цепей переменного тока:
\begin{itemize}
    \item контроль повышенного значения (небаланса) дифференциального тока (является частью функционального блока ДЗТ);
    \item контроль симметрии подводимых к устройству вторичных цепей трансформаторов тока по сторонам защищаемого трансформатора (является отдельным функциональным блоком КЦТ).
\end{itemize}
\item
Параметр <<Номинальный ток токового входа терминала>> – должен соответствовать номинальному току аналоговых входов переменного тока рассматриваемой стороны.
\item
Параметр <<Минимальная величина максимального из фазных токов (LIсим)>> – должен в общем соответствовать току стороны, соответствующему 10\% нагрузки от полной мощности трансформатора.
Остальные значения параметров функции можно оставить заданными по умолчанию производителем.

\end{enumerate}
\end{enumerate}